\runninghead{Passage-by-Passage Exposition}
\section{Detailed Treatment of the Twelve Appointed Parts}

Each appointed part of the First Sunday of Advent in Year~B belongs to its own literary and liturgical setting before it stands alongside the rest of the Mass. The Missal prayers and antiphons serve every First Sunday of Advent across Years~A, B, and~C, whereas the readings, Responsorial Psalm, and Gospel Acclamation are specific to Year~B under the \work{Ordo lectionum Missae}. Each passage raises particular textual questions on which Fathers and Doctors of the Church give distinct answers. Historical settings, traditional attributions, and composition ranges for the scriptural passages are gathered in the \hyperref[sec:date-location]{Scriptural Date and Location} appendix.

\subsection{Entrance Antiphon: \emph{Ad te levavi} (Ps.~25[24]:1--3)}

Psalm~25 (Vulgate~24) opens with an upward turn of desire: \emph{Ad te, Domine, levavi animam meam} (``To thee, O Lord, have I lifted up my soul''), and immediately asks not to be put to shame and not to let enemies mock him, grounding that petition in communal trust: \emph{etenim universi qui sustinent te, non confundentur} (``for none of them that wait on thee shall be confounded''). In Latin, \emph{sustinere} carries the weight of holding up under pressure, enduring while help has not yet arrived.

St.~Augustine, in his \work{Enarrationes in Psalmos} on Psalm~24 (§§1--3), hears Christ speaking here in the person of the Church, lifting up a soul that had been ``trodden down on the earth with carnal longings.'' Augustine reads the fear of shame through the fall of Adam: by trusting in himself and desiring to be as God, man fell into such fragility that he fears death ``from the smallest insect,'' while demonic enemies mock him with secret suggestions and flattering praise. Turning trust away from oneself to God alone reverses that humiliation. Later in the same exposition (§10), commenting on verse~10 --- ``All the ways of the Lord are mercy and truth'' --- Augustine opens a horizon that covers the entire psalm (which also appears among the closed U.S.\ Advent Common Psalms under GIRM~61): ``all the ways of the Lord'' are the two Advents of the Son of God, ``the one in mercy, the other in judgment,'' the first forgiving sins without human desert and the second testing works in incorruptible truth.

St.~Robert Bellarmine (\work{Explanatio in Psalmos}, Ps.~24, explanations~1--3) reads the opening verses along a closely related pastoral line. Finding ``briers and thorns'' everywhere in creatures, the believer tears his soul free from earthly affections and clings to God in temptation. Bellarmine notes that ``my enemies'' encompasses both wicked human adversaries and evil spirits, whose derision would bring intolerable confusion if trust in God proved empty; yet long experience, ancestral examples, and God's own promises assure those who ``wait on the Lord'' --- expecting him in certain hope of assistance --- that they will not be disappointed.

\subsection{Collect: \emph{Da, quaesumus, omnipotens Deus}}

Because the \work{Gloria in excelsis} is omitted on Sundays of Advent (GIRM~53), the Mass passes directly from the Penitential Act (or optional sprinkling, GIRM~51) into the Collect. The modern Latin and English texts of \emph{Da, quaesumus, omnipotens Deus} (``Grant, we beseech thee, almighty God'') are withheld under rights limits and are neither quoted nor paraphrased. For historical background, the ancient Roman \work{Sacramentarium Gelasianum Vetus} (Book~II, lxxxiv, \work{Postcommunio}~683) preserves an early prayer beginning with the same opening words. That prayer asks almighty God to grant his whole family the will (\emph{voluntatem}) to go forth fittingly in righteous deeds (\emph{in operibus iustis aptos occurrere}) to meet Christ his Son at his coming, so that, joined to his right hand, they may possess the heavenly kingdom. Two notes in that historical formula bear directly on this Sunday: even the willingness to meet Christ is petitioned from God as a gift, and the verb \emph{occurrere} (``to run forth to meet'') paired with \emph{in operibus iustis} (``in righteous deeds'') verbally recalls Isaiah~64:5~Vg.\ (Heb.~64:4; \emph{Occurristi laetanti, et facienti justitiam}, ``Thou hast met him that rejoiceth, and doth justice''). The Gelasian formula serves strictly as historical background rather than a substitute for the modern Missal prayer.

\subsection{First Reading: Isaiah~63:16b--17, 19b; 64:2b--7 (Vg.~63:16b--17; 64:1, 3b--8)}

The First Reading draws from the communal lament of Isaiah~63:7--64:12~Vg.\ (Heb.~63:7--64:11), voiced by a people grieving over a ruined sanctuary and an absent God (see \hyperref[sec:date-location]{Scriptural Date and Location}). The Lectionary omits three blocks: 63:16a (the comparison in which Abraham does not know Israel and Israel does not acknowledge them), 63:18--19a (Vg.~63:18--19, lamenting enemies who trampled the sanctuary), and 64:1--2a (Vg.~64:2--3a, the simile of fire kindling brushwood and making water boil, together with the opening clause of Vulgate~64:3, ``When thou shalt do wonderful things, we shall not bear them''). Four acute exegetical difficulties remain, and St.~Jerome (\work{Commentarii in Isaiam} XVII, PL~24 cols.~642--651) and St.~Thomas Aquinas (\work{Expositio super Isaiam ad litteram}, ch.~63--64) resolve three of them in noticeably different ways.

\subsubsection{Divine hardening and wandering (Isa.~63:17)}
Isaiah's first question is daring: \emph{Quare errare nos fecisti, Domine, de viis tuis; indurasti cor nostrum ne timeremus te?} (``Why hast thou made us to err, O Lord, from thy ways: why hast thou hardened our heart, that we should not fear thee?''). Neither Jerome nor Aquinas allows the verse to make God the author of sin, yet they account for the cause differently. Jerome attributes the experience of wandering and hardness to God's \emph{nimia patientia} (``excessive forbearance''): because God awaits our salvation and holds back immediate punishment rather than striking sinners at once, his patience looks to human eyes like a cause of error and hardness --- like a physician hiding the sharp scalpel or a father withholding the rod (citing Hos.~4:14, Hos.~2:14, Ps.~89[88]:33, and Heb.~12). Aquinas, citing Romans~9:18 (``He hath mercy on whom he will, and whom he will he hardeneth''), answers through permissive causality: God makes Israel err and hardens the heart \emph{permittendo, non immittens gratiam, et non corrigens} (``by permitting, not infusing grace, and not correcting''). St.~Augustine (\work{De gratia et libero arbitrio}, ch.~45 [XXIII]) states the doctrinal principle behind both answers on biblical hardening (Exodus~4--14 and Romans~9:18): when God withholds assisting grace in just judgment, free will is not abolished, for ``both God hardened him by His just judgment, and Pharaoh by his own free will.''

\subsubsection{The rending of the heavens and the sudden past tense (Isa.~64:1 $\to$ 64:3b Vg.; 63:19b $\to$ 64:2b Heb.)}
The Lectionary places \emph{Utinam dirumperes caelos, et descenderes} (``O that thou wouldst rend the heavens, and wouldst come down,'' Vg.~64:1) directly beside the past-tense verb \emph{descendisti, et a facie tua montes defluxerunt} (``thou didst come down, and at thy presence the mountains melted away,'' Vg.~64:3b). Jerome and Aquinas diverge directly on what that past tense means. Jerome admires the Hebrew text's sudden shift from wish to accomplished fact as a prayer heard on the spot: no sooner do the speakers cry \emph{Utinam dirumperes} than they address God directly in the past tense --- \emph{Descendisti, Verbum caro factum es, et habitasti in nobis, vere Emmanuel} (``Thou didst come down; thou, the Word, wast made flesh and didst dwell among us, truly Emmanuel''). For Jerome, the ``mountains'' that melt like wax before a consuming fire (Deut.~4:24; Ps.~97[96]:5) are the hostile powers and every proud height raised against the knowledge of God. Aquinas recognizes that Christological sense --- on 64:1~Vg.\ he reads \emph{dirumperes} of Christ's self-emptying (\emph{exinanires temetipsum, deponens majestatem, et carnem assumeres}, ``emptying thyself, laying aside majesty, and taking flesh''), and on 64:3b~Vg.\ he mentions both Final Judgment (\emph{ad judicandum nos}) and the descent into the Virgin's womb (\emph{in uterum virginis}). Yet Aquinas explicitly prefers a historical literal sense (\emph{vel melius, secundum litteralem sensum}) for 64:3b~Vg. Having already linked Israel's dread in the omitted clause 64:3a~Vg.\ (Heb.~64:2a; \emph{non sustinebimus}) to Sinai (Exod.~20), he reads \emph{Descendisti} in 64:3b~Vg.\ as proof drawn from Israel's saving past: if God were to come down now, the same wonders would happen that already happened when he came down to deliver Israel from Egypt (Exod.~4; Ps.~18[17]:10) and melted the cliffs of Arnon (Num.~21).

\subsubsection{Unseen gifts and St.~Paul's paraphrase (Isa.~64:4 Vg.; 64:3 Heb.)}
The verse \emph{A saeculo non audierunt... oculus non vidit, Deus, absque te, quae praeparasti exspectantibus te} (``From the beginning of the world they have not heard... the eye hath not seen, O God, besides thee, what things thou hast prepared for them that wait for thee'') stands behind Paul's words in 1~Corinthians~2:9. Jerome defends Paul's citation as a legitimate sense-for-sense rendering from authentic Scripture (\emph{non verbum ex verbo reddens... sed sensuum exprimens veritatem}, ``not rendering word for word, but expressing the truth of the meaning''), explicitly rejecting apocryphal books such as the \work{Ascensio Isaiae} and \work{Apocalypsis Eliae} as well as Gnostic fabrications (naming Basilides and appealing to St.~Irenaeus, \work{Adversus Haereses} Book~I). Aquinas, quoting 1~Corinthians~2:9 alongside Isaiah~64:4~Vg.\ (Heb.~64:3), specifies four gifts God prepares for those who wait for him: a place of eternal rest (John~14), the kingdom of eternal dignity (Matt.~25), the table of divine refreshment (Ps.~23[22]), and the lamp of eternal light (Ps.~132[131]).

\subsubsection{Inverted anger, stained righteousness, hidden face, and the Potter's clay (Isa.~64:5--8 Vg.; 64:4--7 Heb.)}
Both Jerome and Aquinas mark the clause \emph{Ecce tu iratus es, et peccavimus} (``behold thou art angry, and we have sinned,'' Vg.~64:5) as \emph{ordo praeposterus} (inverted order): we did not sin because God was angry, but God was angry because we sinned. Aquinas adds a brief mystical note on the preceding clause, \emph{Occurristi laetanti, et facienti justitiam} (``Thou hast met him that rejoiceth, and doth justice''), seeing Simeon in the Temple (Luke~2). On verse~6~Vg.\ (Heb.~v.~5; \emph{quasi pannus menstruatae universae justitiae nostrae}, ``all our justices as the rag of a menstruous woman''), the two Doctors divide again. Jerome reads the harsh simile against legal observance after the coming of Christ: compared with evangelical purity, merely legal righteousness is counted unclean, as St.~Paul counted his legal attainments loss for Christ (Phil.~3:8) and as Esther loathed her royal diadem (Esther~14:16). Aquinas instead examines how any genuinely good act of justice can become stained in four concrete ways: by mortal sin (Sir.~47), by negligence (Jer.~48), by a twisted intention (Matt.~5), or by rising pride (Isa.~57).

In verse~7~Vg.\ (Heb.~v.~6), the prophet reaches the lowest point: \emph{Non est qui invocet nomen tuum; qui consurgat, et teneat te. Abscondisti faciem tuam a nobis} (``There is none that calleth upon thy name: that riseth up, and taketh hold of thee: thou hast hid thy face from us''). Commenting on this exact verse (PL~24 col.~650), Jerome exhorts sinners to confess their iniquity and immediately quotes Psalm~80(79):4 --- \emph{Ostende faciem tuam, et salvi erimus} (``Show thy face, and we shall be saved'') --- linking Isaiah's lament directly to the verse that serves as this Sunday's Responsorial Psalm refrain. Finally, verse~8~Vg.\ (Heb.~v.~7) turns the lament toward hope without appealing to human merit: \emph{nos vero lutum; et fictor noster tu, et opera manuum tuarum omnes nos} (``and we are clay: and thou art our maker, and we all are the works of thy hands''). Jerome observes that examining our own deserts leads to despair, whereas calling on God's Fatherhood (Exod.~4:22) and Potter's craft leaves room for prayer; Aquinas pairs the verse with Jeremiah~18 (``like clay in the potter's hand'') and Job~10.

\subsection{Responsorial Psalm: Psalm~80(79):2ac, 3b, 15--16, 18--19; R.~v.~4}

Psalm~80 (Vulgate~79) is a communal plea for Israel's restoration, structured around its recurring refrain: \emph{Deus, converte nos, et ostende faciem tuam, et salvi erimus} (``Convert us, O God: and shew us thy face, and we shall be saved''). The Year~B Lectionary trims the psalm sharply, omitting the names of Joseph, Ephraim, Benjamin, and Manasseh (vv.~2b, 3a), the middle section depicting the bread of tears, mocking neighbors, and the vine brought out of Egypt and torn apart by the wild boar (vv.~5--14), and the imprecation of verse~17. What remains moves directly from invoking the Shepherd enthroned upon the Cherubim to asking him to visit his vineyard and place his hand upon the ``son of man'' (\emph{filium hominis}).

St.~Augustine (\work{Enarrationes in Psalmos}, Ps.~79, §§1--3, 9, 11) reads the entire psalm as a prophecy of Christ and the Church: ``The song here is of the Advent of the Lord and of our Saviour Jesus Christ, and of His vineyard... Head and Body, King and people, Shepherd and flock.'' On verse~2, Augustine notes that \emph{Cherubim} signifies ``fulness of knowledge,'' which is not restricted to angelic powers: because ``The fulness of the Law is love'' (Rom.~13:10), any believer who holds fast to charity becomes God's throne. On verse~3b (\emph{Excita potentiam tuam}, ``Stir up thy might''), Augustine thinks of the Cross, where Christ seemed weak like the angel who let Jacob prevail in wrestling so as to bless him, before rising in power. On the refrain (v.~4, ``illumine Thy face''), he asks how God's face could ever be darkened, and answers that Christ placed before it ``a cloud of flesh, and as it were a veil of weakness,'' unrecognized on the tree until the Resurrection and Ascension unveiled who he was. On verses~15--16, Augustine stresses that God plants no other vine, but perfects the vine of Abraham's seed ``upon the Son of man'' --- Christ Jesus, the sole foundation (1~Cor.~3:11). And on verse~19 (\emph{vivificabis nos, et nomen tuum invocabimus}, ``thou shalt quicken us: and we will call upon thy name''), Augustine insists that true life means loving God \emph{gratis} (for himself alone) rather than for earthly goods he gives to wicked and good alike. Where Isaiah~64:7~Vg.\ (Heb.~64:6) laments that no one calls upon God's name, Psalm~80(79):19 ends with the vow that the quickened people will call upon his name again.

St.~Robert Bellarmine (\work{Explanatio in Psalmos}, Ps.~79, explanations~1--3 and 14--18) complements Augustine's ecclesiological exposition with historical and literal-messianic analysis. He grounds the omitted tribal names in the geography of the land of promise (Samaria in Ephraim, Jerusalem in Benjamin, the territory beyond the Jordan in Manasseh) and links ``thou that sittest upon the Cherubim'' to the ark's mercy seat. On verses~15--19, Bellarmine argues that ``All agree that this passage refers literally to Christ,'' because after the Babylonian captivity the Jewish people never recovered their former position, so the vineyard was made perfect only under Christ. He links ``visit this vineyard'' (v.~15) to Zechariah's canticle (``He hath visited and wrought the redemption of his people,'' Luke~1:68) and to Christ weeping over Jerusalem for not knowing ``the time of thy visitation'' (Luke~19:44), and he reads ``the man of thy right hand'' (v.~18) of Christ formed by the right hand of God alone in the Virgin's womb through the Holy Spirit (Luke~1:35).

Where local parishes sing one of the U.S.\ seasonal Advent Common Psalms permitted by GIRM~61 (Psalm~25[24] or Psalm~85[84], under U.S.\ Lectionary No.~174), the assembly prays psalms already represented in this Mass's Entrance Antiphon, Gospel Acclamation, and Communion Antiphon.

\subsection{Second Reading: 1~Corinthians~1:3--9}

St.~Paul's thanksgiving to the church in Corinth stands immediately before his sharp reprimand of their factions in 1:10--12 and their spiritual immaturity in 3:1--3 (see \hyperref[sec:date-location]{Scriptural Date and Location}). That juxtaposition raises two questions about 1~Corinthians~1:3--9: why does Paul thank God that a quarrelsome community has been ``made rich in him, in all utterance and in all knowledge'' (1:5) and lacks no grace (1:7), and what does he mean by promising that faithful God will confirm them ``unto the end without crime'' (1:8--9)?

St.~John Chrysostom (\work{Homilies on First Corinthians}, Homily~II, §§1--8) faces the tension between 1:4--7 and 1:10 directly. He offers three complementary explanations: first, the Corinthians had shown real zeal at the beginning before growing lax; second, Paul praises the faithful remnant among them (such as the household of Stephanas) without naming and shaming the offenders at the opening of the letter; third, Paul uses pastoral management (\emph{oikonomikos}), attributing every single gift to ``the grace of God'' (\emph{in gratia Dei}) to cut down their vanity before correcting them. Chrysostom points out how Paul ``fasteneth them as with nails to the Name of Christ'' across verses~1--10, setting himself ``as it were to arouse those who were heavy-headed after some debauch,'' and applying that Name incessantly so that he may ``foment their inflammation, and purge out the corruption of the disease'' (§7) --- curing their habit of naming themselves after Paul or Apollos.

On verse~7 --- \emph{exspectantibus revelationem Domini nostri Jesu Christi} (``waiting for the manifestation of our Lord Jesus Christ'') --- Chrysostom fastens on the exact noun \emph{revelation} (Greek \emph{apokalupsis}, ``unveiling''): ``implying that although He be not seen, yet He is, and is present even now, and then shall appear.'' Waiting for Christ is waiting for a veil to be drawn back from a Lord who is already here. On verse~8, Chrysostom notes that promising to ``confirm'' them hints quietly that they are still wavering and in danger; and on verse~9 (\emph{Fidelis Deus}, ``God is faithful''), he rejects any fatalistic reading of perseverance: God is faithful to bring believers to the promised goods ``if we be not quite impatient of His rein,'' echoing Paul's own words from 1~Corinthians~2:9 (``which eye hath not seen, nor ear heard,'' in Paul's eye-before-ear order, recalling Isa.~64:4~Vg.\ [Heb.~64:3]).

Cornelius a Lapide (\work{Commentaria in omnes divi Pauli epistolas}, Vol.~7, on 1~Cor.~1:3--9) reads Paul's praise on verses~4--5 more directly of the genuine spiritual wealth distributed across the Corinthian community as a whole, softening Chrysostom's irony. On verses~8--9, however, Lapide agrees with Chrysostom in protecting human cooperation by citing the Council of Trent (Session~VI, chapters~11 and~13, following Augustine): God confirms believers ``so far as His part is'' and will not desert them unless he is first deserted, so the promise forbids despair without granting presumptuous certainty apart from special revelation.

\subsection{Gospel Acclamation: \emph{Ostende nobis, Domine} (Ps.~85[84]:8)}

Outside Lent, the acclamation before the Gospel retains the \emph{Alleluia} (GIRM~62). Its verse, \emph{Ostende nobis, Domine, misericordiam tuam, et salutare tuum da nobis} (``Shew us, O Lord, thy mercy; and grant us thy salvation''), repeats the verb \emph{ostende} (``show'') from the Responsorial Psalm refrain (\emph{ostende faciem tuam}), asking God to reveal his mercy and give his saving gift.

St.~Augustine (\work{Enarrationes in Psalmos}, Ps.~84, §6) interprets ``thy salvation'' without hesitation: ``Thy salvation, that is, Thy Christ.'' Yet Augustine immediately turns the prayer inward against pride. Happy is the person to whom God shows his mercy, because once a man sees that every good thing praised in him comes from mercy rather than desert, he is not puffed up; not being lifted up, he does not fall; standing fast, he clings to God. Faith now purifies the heart so that when the light of Christ is shown, it may find us whole --- seeing Christ ``not as the Jews saw Him and crucified Him, but as the Angels see Him, and rejoice.'' St.~Robert Bellarmine (\work{Explanatio in Psalmos}, Ps.~84, explanation~7) places the verse squarely in the expectation of the Incarnation: asking God to ``show'' his mercy asks him to make his eternal love visibly experienced by sending his only-begotten Son (citing Titus~2:11 and~3:4), while commending Augustine's moral sense as its necessary companion.

\subsection{Gospel: Mark~13:33--37 (with omitted Mark~13:32)}

Mark~13:33--37 closes Jesus' Olivet Discourse spoken opposite the Temple (see \hyperref[sec:date-location]{Scriptural Date and Location}). The Lectionary opens at verse~33, omitting verse~32: \emph{De die autem illo vel hora nemo scit, neque angeli in caelo, neque Filius, nisi Pater} (``But of that day or hour no man knoweth, neither the angels in heaven, nor the Son, but the Father''). Patristic and medieval commentators read verses~33--37 continuously out of verse~32, confronting three linked questions: how the Son can speak of not knowing the hour, in what sense the householder has ``gone abroad'' and left his house, and how the closing charge --- ``And what I say to you, I say to all: Watch'' --- applies to believers across every century.

\subsubsection{The withheld hour as an economy of vigilance (Mark~13:32--33)}
In Book~IX of \work{De Trinitate} (§§58--68), St.~Hilary of Poitiers answers Arian arguments that inferred an imperfect divine nature from Christ's words about the unknown day and hour (Mark~13:32; Matt.~24:36). Hilary asks how the Lord of glory, in whom all treasures of wisdom and knowledge are hidden (Col.~2:3), could be ignorant of the day of his own coming. When Scripture speaks of God ``not knowing,'' Hilary argues, it marks an economy (\emph{dispensatio}) of silence or withholding action rather than a defect of nature. Pairing that silence with Christ's Matthean watchfulness sayings --- ``Watch therefore, for ye know not the day nor the hour'' (Matt.~25:13, cited in IX.65), ``Therefore be ye also ready'' (Matt.~24:44), and ``Blessed is that servant'' (Matt.~24:46, both in IX.67) --- Hilary explains in IX.67 that Christ ``withholds from us the security of certain knowledge, that our minds may be kept on the stretch by the uncertainty of suspense... for the security of knowledge might breed negligence of the faith.'' Centuries later, St.~Thomas Aquinas in the \work{Catena Aurea in Marcum} (13:32--37) applies Hilary's reasoning directly to Mark~13:33 (``Take ye heed, watch and pray'').

\subsubsection{Bodily departure, spiritual authority, and the doorkeeper (Mark~13:34--36)}
St.~Bede the Venerable, in \work{In Marci Evangelium Expositio} Book~IV, chapter~13, adopts Hilary's anti-Arian line, observing that Christ says \emph{nescitis} (``you do not know''), not ``we do not know,'' so that disciples remain in \emph{pendulae exspectationis incerto} (``the uncertainty of poised expectation''). Turning to the parable of the \emph{homo qui peregre profectus reliquit domum suam} (``a man who, going into a far country, left his house''), Bede clarifies the departure by borrowing directly from St.~Gregory the Great's Homily~9 on Matthew~25:14--30 (\work{Homiliae in Evangelia}, Hom.~9): earth is the proper home of human flesh, which was carried as if to a far country when the risen Redeemer exalted our flesh to heaven at the Ascension. Christ therefore left his Church only bodily (\emph{corporaliter}), while never withdrawing the safeguard of his divine presence.

Within that household, Bede reads \emph{potestatem cujusque operis} (``authority over every work'') as the grace of the Holy Spirit empowering the faithful for every good work, and identifies the \emph{janitor} (``porter'' or doorkeeper) with the order of pastors charged to guard the flock. Yet verse~37 broadens that office to every believer guarding the door of the heart against evil suggestions. To define what watching looks like, Bede quotes St.~Gregory the Great's Homily~13 (on the Lucan parallel Luke~12:35--40): \emph{Vigilat qui ad aspectum veri luminis mentis oculos apertos tenet, vigilat qui servat operando quod credit, vigilat qui a se torporis et negligentiae tenebras repellit} (``He watches who keeps the eyes of his mind open to the sight of the true light; he watches who keeps in deeds what he believes; he watches who drives away from himself the darkness of torpor and negligence'').

\subsubsection{The four watches and the command ``to all'' (Mark~13:35--37)}
St.~Gregory the Great, Theophylact of Ohrid, Pseudo-Jerome, and St.~Augustine each address a further question raised by the closing verses. First, St.~Gregory the Great (\work{Homiliae in Evangelia}, Hom.~13) reads the watches of the night as the stages of human life --- boyhood, youth, and old age --- showing that God is a \emph{patiens redditor} (``patient judge'') who leaves no age without hope of conversion yet demands stricter account if his patience is spurned, keeping the final hour unknown \emph{ut semper possit esse suspecta} (``so that it may always be anticipated''). In Homily~1 (on Luke~21:25--33, §§1 and~5), Gregory reinforces the same detachment by portraying the world in its old age, worn down with growing afflictions so that believers stop clinging to what cannot last.

Second, St.~Thomas Aquinas's \work{Catena Aurea in Marcum} (13:32--37) transmits two additional early commentators on verses~35--37. Theophylact of Ohrid warns that many stay awake only to pass the night in wickedness --- which is why Christ joins watching with prayer --- and distributes the four Roman night watches across human life in a different order from Gregory: evening is old age, midnight is mid-youth, cockcrow is the awakening of mature reason, and morning is childhood (showing that even a child needs watchful care). Pseudo-Jerome adds a psychological image of sleep: a sleeper gives his mind to phantoms that vanish upon waking, just as those seized by worldly love lose after death everything they dreamed was real.

Third, through the same section of the \work{Catena Aurea}, Aquinas preserves St.~Augustine's answer in \work{Epistula}~199 §3 to a sharp objection: how can Jesus say ``Watch, lest when he cometh he find you sleeping?'' to \emph{all} believers of every century when only one generation will be alive at the Last Day? Augustine answers that ``that day comes to each man when his day comes for departing from this life such as he is to be, when judged in that day.'' Because death fixes the state in which each person will meet the Judge, every generation stands equally under the command to watch.

\subsection{Prayer over the Offerings: \emph{Suscipe, quaesumus, Domine}}

The modern Latin and English texts of \emph{Suscipe, quaesumus, Domine} (``Accept, we beseech thee, O Lord'') are withheld under rights limits, and no proper Offertory antiphon is printed in the Missal (GIRM~74 leaves the Offertory chant to local selection under the norms for the Entrance chant). The ancient Roman formula of the \work{Sacramentarium Veronense} (Section~XXVIIII, marginal~372) preserves its historical antecedent, differing in its petition body from the secondary 2002 digital reproduction of the \work{Missale Romanum} (pp.~82--83) only by a single inflection --- ancient \emph{temporalis} against secondary-2002 \emph{temporali} --- while correspondence with the printed 2002 and printed 2008 editions remains unverified. In that ancient prayer, the assembly offers gifts \emph{quae de tuis offerimus collata beneficiis} (``which we offer from thy bestowed benefits'') and asks that what God grants to be accomplished for our temporal devotion may become the reward of eternal redemption. Even at the altar, nothing human is brought to God that did not first come from his generosity, and temporal worship is ordered toward eternal life.

\subsection{Seasonal Preface: The Twofold Horizon of Advent}

The \work{Ordo Missae} provides two Prefaces of Advent, and GIRM~364 notes that the Missal's Prefaces bring out the motives for thanksgiving and the facets of the mystery of salvation. Because the modern Preface text is withheld under rights limits, three patristic and medieval sermons on Christ's comings illuminate Advent thanksgiving directly.

St.~Cyril of Jerusalem (\work{Catechesis}~15, §§1--2, 4, 26) teaches Christ's two Advents: ``We preach not one advent only of Christ, but a second also, far more glorious than the former.'' Almost everything in Christ is twofold --- a twofold generation (of God before the ages, of a Virgin at the close of the ages) and twofold descents: ``one, the unobserved, like rain on a fleece; and a second His open coming, which is to be.'' In the first he lay wrapped in swaddling clothes in the manger and endured the Cross; in the second he comes clothed in light, attended by a host of angels, bringing the crown of a divine kingdom (§1), as Malachi~3:1--2 foretold (§2). Cyril forbids curiosity about dates from Acts~1:7 (§4), and in §26 insists from Matthew~25:35--41 that ``there is here no need of allegory, but to fulfil what is said'' through works of mercy, keeping one's lamp constantly burning.

In the Roman Church, St.~Leo the Great (\work{Sermon}~19, §§1--3, preached for the December fast) brings that two-Advent watchfulness into parish life: quoting Luke~21:34 against hearts weighed down by gluttony and worldly cares because ``the day denounced is near, even though hidden'' (§1), he urges believers to pair bodily fasting with a fast of the mind from vice (§2) and generous almsgiving to the poor (§3). St.~Bernard of Clairvaux (\work{Sermones in Adventu}~3.4 and~5.1) describes a middle coming between the Incarnation and the Final Judgment. In Sermo~3.4, Bernard distinguishes three comings of the Lord --- \emph{ad homines, in homines, contra homines} (``to men, into men, against men'') --- and quotes John~14:23 to show how Christ comes spiritually and invisibly into those who love him and keep his word. In Sermo~5.1, expounding the psalmist's promise to ``sleep among the midst of lots'' (Ps.~68[67]:14), he contrasts the first coming ``in the flesh and in weakness'' and the last ``in glory and majesty'' with this hidden middle advent \emph{in spiritu et virtute} (``in spirit and in virtue,'' or power), visible only to the elect as the road from redemption to final glory.

\subsection{Communion Antiphon: \emph{Dominus dabit benignitatem} (Ps.~85[84]:13)}

The Communion Antiphon --- \emph{Dominus dabit benignitatem, et terra nostra dabit fructum suum} (``The Lord will give goodness, and our earth shall yield her fruit,'' Ps.~85[84]:13) --- comes from the same Korahite psalm as the Gospel Acclamation. St.~Augustine and St.~Robert Bellarmine expound the verse with different emphases.

St.~Augustine (\work{Enarrationes in Psalmos}, Ps.~84, §§10--12) reads the preceding verse --- ``Truth hath sprung out of the earth'' (v.~12) --- in two ways: first Christologically, of the Son of God taking mortal flesh from the Virgin Mary so that he might have a body to offer as a healing sacrifice on the Cross (§10); and second morally, of sinful human ``earth'' speaking truth by confessing its sins so that God's righteousness looks down in pardon (§11). Then on verse~13 (§12), rendering \emph{benignitas} as ``sweetness,'' Augustine explains that God gives the interior delight of working righteousness, transforming tastes so that one who once delighted in drunkenness, theft, or idle shows now delights in sobriety, giving, prayer, and hymns. Only when God rains that grace upon hearts ``furrowed, as it were, with the plough of confession'' does ``our earth'' yield her harvest.

St.~Robert Bellarmine (\work{Explanatio in Psalmos}, Ps.~84, explanation~12) keeps the verse tightly anchored in the mystery of the Incarnation: ``our earth'' is the untouched Virgin Mary, never ploughed or sown by human seed, whom the Holy Spirit overshadows from heaven so that she brings forth Christ spontaneously like the ``flower of the field'' (Song~2:1). Placed at Holy Communion, both expositions illuminate the sacrament without competing: Christ, the fruit born of Mary, enters the soil of the communicant's heart so that grace may make righteousness sweet and fruitful.

\subsection{Prayer after Communion (\emph{Prosint nobis}) and Concluding Blessing (GIRM~90)}

The Prayer after Communion (\emph{Prosint nobis}, ``May they profit us'') closes the Communion Rite (\work{Ordo Missae} nn.~124--146). Because its modern wording is withheld under rights limits and no ancient Roman sacramentary ancestor has been established, no specific petitions are attributed to it beyond its opening words and its place after Holy Communion.

The Concluding Rites then follow under GIRM~90: brief announcements if necessary, the Priest's greeting and blessing --- which on particular days and occasions may be enriched by an optional solemn blessing or prayer over the people (introduced under the 2008 \work{Notitiae} rubric by \emph{Inclinate vos ad benedictionem}, ``Bow down for the blessing'') --- and the dismissal (with the three alternative formulas added after \emph{Ite, missa est} in 2008), sending each member forth to live in good works and praise of God. The modern Advent solemn-blessing wording is withheld under rights limits and is neither quoted nor paraphrased.
